\documentclass[12pt]{article}

\usepackage{sbc-template}

\usepackage{graphicx,url}
\usepackage[utf8]{inputenc}
\usepackage[brazil]{babel}
\usepackage{amsmath}
\usepackage{booktabs}
\usepackage{hyperref}

\sloppy

\title{Previsão de Crimes Violentos e de Sinistros de Trânsito em Dias
de Jogos de Futebol: Proposta de um Sistema de Ciência de Dados para a
Cidade de São Paulo}

\author{Gilberto Serafim Dias de Araújo\inst{1}, Leonardo de Cico\inst{1},
Pedro Henrique de Almeida\inst{1}}

\address{Universidade de São Paulo (USP)\\
  São Paulo -- SP -- Brasil
  \email{serafimgilberto0@usp.br, leocico@usp.br, pedroalmeida04@usp.br}
}

\begin{document}

\maketitle

\begin{resumo}
Partidas de futebol concentram multidões em torno de estádios e áreas de
entretenimento e alteram, por algumas horas, as rotinas de deslocamento e
convívio de uma cidade inteira. A literatura em criminologia ambiental mostra
que essa alteração de rotina está associada a variações mensuráveis na
incidência de crimes e de sinistros de trânsito. Este trabalho propõe a
arquitetura de um sistema preditivo que integra, em uma tabela analítica de
granularidade diária, o calendário de partidas dos clubes da cidade de São
Paulo, os registros oficiais de ocorrências criminais e de sinistros de
trânsito, e variáveis de clima e de calendário. São descritas as fontes de
dados públicas, o pipeline de engenharia de dados, a formulação do problema
preditivo e o protocolo de comparação entre um modelo baseado em árvores de
decisão (XGBoost) e um modelo aditivo de séries temporais (Prophet). A
proposta adapta ao contexto brasileiro o desenho experimental de estudos
internacionais que relacionam eventos esportivos e criminalidade espacial.
\end{resumo}

\begin{abstract}
Football matches concentrate crowds around stadiums and entertainment areas
and disrupt, for a few hours, the commuting and social routines of an entire
city. Environmental criminology literature shows that this routine disruption
is associated with measurable variation in crime and traffic accident
incidence. This work proposes the architecture of a predictive system that
integrates, into a single daily analytical table, the match calendar of São
Paulo football clubs, official records of criminal occurrences and traffic
accidents, and weather and calendar variables. We describe the public data
sources, the data engineering pipeline, the formulation of the predictive
problem and the protocol for comparing a tree-based model (XGBoost) with an
additive time-series model (Prophet). The proposal adapts to the Brazilian
context the experimental design of international studies relating sporting
events and spatial crime analysis.
\end{abstract}

\section{Introdução}

Grandes eventos esportivos são, do ponto de vista urbano, deslocamentos
temporários de rotina. Em poucas horas, dezenas de milhares de pessoas se
concentram no entorno de um estádio, alteram trajetos habituais, consomem
álcool em bares e praças de concentração e retornam simultaneamente para casa
ao final da partida. A Teoria da Atividade Rotineira \cite{cohen1979} descreve
exatamente esse cenário como propício à ocorrência de crimes: a convergência,
no tempo e no espaço, de ofensores motivados, alvos disponíveis e ausência de
vigilância capaz. São Paulo é um caso extremo desse fenômeno no Brasil, por
concentrar quatro clubes de série A, quatro estádios de grande porte e uma
agenda que, entre campeonatos estadual, nacional, copas e competições
continentais, raramente deixa uma semana sem jogo na cidade. Ao mesmo tempo, a
disponibilidade de dados abertos sobre segurança pública, sinistros de trânsito
e clima tornou tecnicamente viável investigar essa relação com métodos
quantitativos, o que até pouco tempo dependia de coleta manual.

O problema tratado neste trabalho é a antecipação, em nível diário e para o
município de São Paulo, de dias atípicos de violência e de sinistros de
trânsito associados à realização de partidas de futebol. A pergunta que
orienta a investigação é concreta: conhecendo o calendário de jogos, as
características de cada partida (mandante, horário, competição, rivalidade,
público esperado), as condições meteorológicas previstas e o calendário civil,
é possível estimar com antecedência de um dia se determinada data apresentará
volume de ocorrências acima do padrão? A delimitação é deliberada. Não se
pretende explicar as causas da criminalidade urbana nem localizar
espacialmente cada ocorrência; a unidade de análise é o dia na cidade, e o
produto é uma estimativa de risco utilizável no planejamento de efetivo
policial e de fiscalização de trânsito.

Três dificuldades justificam a pesquisa. A primeira é metodológica: os
trabalhos de referência sobre o tema \cite{kurland2014,ristea2018,ristea2020}
são espaciais e explicativos, construídos sobre células de grade no entorno do
estádio e voltados a medir a associação entre evento e crime, não a produzir
previsões operacionais para a cidade como um todo. A segunda é de
disponibilidade de dados: parte relevante do ganho preditivo relatado nesses
estudos vem de mensagens geolocalizadas de redes sociais, fonte cujo acesso
programático foi drasticamente restringido nos últimos anos, o que torna
necessário verificar quanto do sinal permanece recuperável apenas com dados
oficiais e de calendário. A terceira é a escassez de estudos que tratem crimes
e sinistros de trânsito como alvos de um mesmo pipeline: no contexto de dias de
jogo, ambos compartilham os mesmos determinantes -- aglomeração, consumo de
álcool, deslocamento noturno concentrado --, mas costumam ser estudados por
comunidades separadas e sobre bases não integradas.

O objetivo geral é propor e especificar um sistema de Ciência de Dados que
estime, em granularidade diária, o risco de picos de crimes violentos e de
sinistros de trânsito na cidade de São Paulo a partir do calendário de partidas
de futebol e de variáveis contextuais. Desdobra-se em quatro objetivos
específicos: (i) levantar e caracterizar as fontes públicas necessárias,
verificando a granularidade efetivamente disponível em cada uma; (ii)
especificar um pipeline de integração que consolide fontes de granularidades
distintas em uma tabela analítica diária única; (iii) formular a tarefa
preditiva e o critério de definição de pico, com protocolo de validação
temporal; e (iv) definir o desenho de comparação entre um modelo tabular
baseado em árvores e um modelo aditivo de séries temporais, de modo a isolar a
contribuição das variáveis de jogo em relação a um modelo apenas sazonal.

O restante do trabalho está organizado da seguinte forma. A Seção 2 revisa os
trabalhos relacionados que fundamentam o desenho da proposta e explicita como
esta se posiciona em relação a eles. A Seção 3 apresenta a metodologia, com as
fontes de dados, o pipeline de engenharia de dados, a formulação do problema
preditivo e o protocolo de avaliação. A Seção 4 apresenta o cronograma de
execução previsto.

\section{Trabalhos Relacionados}

A relação entre eventos esportivos e criminalidade urbana já foi estudada sob
perspectivas metodológicas distintas. Kurland et al. \cite{kurland2014}
examinaram o entorno do Villa Park (Aston Villa Football Club, Birmingham) e
mostraram que a distribuição espacial de furtos e de crimes violentos contra a
pessoa difere de forma significativa entre dias de jogo e dias de comparação em
um raio de aproximadamente três quilômetros do estádio. O estudo é relevante
aqui por dois motivos: estabelece o desenho de dias de evento contra dias de
controle equivalentes, adotado nesta proposta, e evidencia que o efeito do jogo
não se restringe ao perímetro imediato do estádio, o que sustenta a escolha de
uma unidade de análise mais ampla.

Ristea et al. \cite{ristea2018} propuseram um dos primeiros arcabouços que
combinam dados criminais georreferenciados com mensagens do Twitter para
explicar a ocorrência de crimes ao redor de um estádio. Ocorrências policiais e
mensagens geolocalizadas da região do Villa Park foram agregadas em células de
500\,m $\times$ 500\,m, comparando-se oito horas em torno do início da partida
em dias de jogo e em dias de controle. Por meio do índice de Moran bivariado e
de modelos de regressão binomial negativa, os autores mostraram que a densidade
de mensagens de conteúdo violento -- identificadas por dicionário léxico --
apresenta correlação espacial mais forte com dano criminal, furto e violência
contra a pessoa do que o volume total de mensagens. O achado sustenta um
princípio adotado nesta proposta: subconjuntos qualificados de uma fonte
dinâmica (aqui, clássicos e jogos noturnos, e não simplesmente ``houve jogo'')
tendem a ser mais informativos do que o agregado indiferenciado.

Em trabalho posterior, Ristea et al. \cite{ristea2020} estenderam a abordagem
para Chicago, cruzando criminalidade, mensagens geolocalizadas e variáveis
demográficas, socioeconômicas e ambientais no entorno do United Center. Com
estimação de densidade de kernel localizada, construíram modelos de previsão
espacial para sete tipos de crime, comparando dias de jogo em casa, jogos fora
e dias de controle. A inclusão dos dados de redes sociais elevou a AUC de
previsão para quase todos os tipos de crime, com ganhos maiores para furto de
veículos e roubo e ganhos reduzidos para crimes já fortemente concentrados no
espaço e no tempo. Essa heterogeneidade motiva, nesta proposta, a análise
desagregada por categoria de ocorrência: é esperado que o efeito das variáveis
de jogo seja desigual entre crimes patrimoniais em via pública e crimes de
outra natureza.

Os três trabalhos convergem em limitações que também condicionam esta proposta:
a defasagem e a granularidade heterogênea entre bases, e o problema da unidade
de área modificável (MAUP), pelo qual a escolha da unidade espacial pode
alterar substancialmente os resultados. O posicionamento deste trabalho em
relação a eles é, portanto, duplo. Metodologicamente, desloca-se do objetivo
explicativo e espacial para um objetivo preditivo e temporal, adotando o dia na
cidade como unidade de análise e a previsão de pico como tarefa. Quanto aos
dados, substitui-se a fonte dinâmica de redes sociais -- indisponível em escala
para o caso brasileiro -- por atributos de partida conhecidos com antecedência,
o que tem a vantagem colateral de tornar o sistema utilizável de forma
prospectiva, e não apenas retrospectiva. Acrescenta-se, por fim, o tratamento
conjunto de crimes e de sinistros de trânsito como alvos do mesmo pipeline.

\section{Metodologia}

\subsection{Fontes de dados}

O sistema integra quatro categorias de dados públicos referentes ao município
de São Paulo:

\begin{itemize}
  \item \textbf{Partidas de futebol}: histórico de jogos do Campeonato
  Brasileiro, do Campeonato Paulista, da Copa do Brasil e das competições
  continentais envolvendo os clubes da cidade, obtido de bases estruturadas
  públicas, com data, horário, mandante, visitante, estádio, competição,
  público e placar;
  \item \textbf{Ocorrências criminais}: registros da Secretaria de Segurança
  Pública do Estado de São Paulo (SSP-SP), com recorte para roubo, furto, lesão
  corporal dolosa e homicídio doloso;
  \item \textbf{Sinistros de trânsito}: registros do Sistema de Informações
  Gerenciais de Sinistros de Trânsito (Infosiga SP), com data, hora, tipo de
  via, localização e gravidade;
  \item \textbf{Clima e calendário}: temperatura, precipitação e umidade
  provenientes de estações meteorológicas públicas, complementadas por dia da
  semana, feriado e véspera de feriado.
\end{itemize}

A primeira atividade do projeto é a verificação da granularidade efetivamente
publicada em cada fonte, porque dela depende a viabilidade do recorte diário. A
base de sinistros e a de partidas oferecem resolução de data e hora; já os
indicadores criminais são frequentemente divulgados em agregações mensais por
município, sendo necessário recorrer às bases de boletins de ocorrência
individualizados. Caso a desagregação diária não se confirme para alguma
categoria de crime, essa categoria é excluída do alvo, e não substituída por
interpolação, procedimento que introduziria variação artificial justamente na
variável a ser prevista.

\subsection{Pipeline de engenharia de dados}

Propõe-se um esquema relacional normalizado, com tabelas independentes para
partidas, ocorrências criminais, sinistros de trânsito e calendário/clima,
consolidadas por consultas SQL em uma tabela analítica única na qual cada linha
representa um dia na cidade de São Paulo. A Tabela~\ref{tab:variaveis} resume as
variáveis previstas.

\begin{table}[ht]
\centering
\caption{Variáveis da tabela analítica diária}
\label{tab:variaveis}
\small
\begin{tabular}{@{}llll@{}}
\toprule
\textbf{Variável} & \textbf{Tipo} & \textbf{Fonte} & \textbf{Função} \\
\midrule
data & Data & Calendário & Chave da tabela analítica \\
houve\_jogo & Binário & Base de partidas & Indicador de evento \\
n\_jogos & Inteiro & Base de partidas & Múltiplas partidas no dia \\
classico & Binário & Base de partidas & Jogo entre grandes rivais \\
publico & Numérico & Base de partidas & Magnitude da aglomeração \\
periodo\_partida & Categórico & Base de partidas & Tarde ou noite \\
competicao & Categórico & Base de partidas & Estadual, nacional, copa \\
resultado\_mandante & Categórico & Base de partidas & Vitória, empate, derrota \\
precipitacao & Numérico & Meteorologia & Inibidor de circulação \\
temperatura\_max & Numérico & Meteorologia & Controle sazonal \\
dia\_semana & Categórico & Calendário & Sazonalidade semanal \\
feriado & Binário & Calendário & Alteração de rotina \\
lag\_1, lag\_7 & Numérico & Alvos & Defasagens do alvo \\
crimes\_dia & Inteiro & SSP-SP & Variável-alvo \\
sinistros\_dia & Inteiro & Infosiga SP & Variável-alvo \\
\bottomrule
\end{tabular}
\end{table}

Duas decisões de projeto merecem registro. A primeira é a construção explícita
de variáveis defasadas ($lag_1$, $lag_7$): modelos tabulares não possuem noção
intrínseca de ordem temporal, e sem essas variáveis o efeito do jogo seria
confundido com o nível corrente da série. A segunda é a definição do conjunto
de dias de controle, formado por dias sem jogo pareados por dia da semana, mês
e condição meteorológica, de modo que a comparação não capture apenas o fato de
que grande parte das partidas ocorre em fins de semana.

\subsection{Formulação do problema preditivo}

Seguindo \cite{ristea2020}, a tarefa é formulada como classificação binária: a
cada dia $d$ atribui-se rótulo $1$ quando o volume de ocorrências supera um
limiar de pico, e $0$ caso contrário. O limiar é definido a partir da
distribuição empírica do próprio alvo -- percentil 90 da série de referência do
mesmo mês --, e não por valor absoluto arbitrário, o que evita que a
sazonalidade anual concentre todos os picos em um único período. A
probabilidade de pico é modelada como função de um vetor de variáveis
explicativas $f_1(d), \ldots, f_n(d)$:

\begin{equation}
P(\text{pico}_d = 1 \mid f_1(d), \ldots, f_n(d)) =
\frac{1}{1 + e^{-(\beta_0 + \sum_{i=1}^{n}\beta_i f_i(d))}}
\end{equation}

A formulação permite comparar subconjuntos de variáveis explicativas: um
modelo apenas com defasagens e calendário, contra um modelo que incorpora as
variáveis de partida. A diferença de desempenho entre os dois é a medida da
contribuição incremental do futebol, de forma análoga à comparação entre
modelos com e sem dados de redes sociais em \cite{ristea2020}. O mesmo conjunto
de variáveis é avaliado também na formulação de regressão sobre a contagem
diária, o que preserva informação descartada pela binarização.

\subsection{Modelos candidatos e protocolo de avaliação}

Comparam-se dois algoritmos de naturezas complementares. O XGBoost, conjunto de
árvores de decisão treinadas sequencialmente por gradiente, é adequado a dados
tabulares heterogêneos e captura interações não lineares sem especificação
manual -- por exemplo, o efeito conjunto de chuva, horário noturno e clássico
sobre o risco de sinistros. O Prophet, modelo de decomposição aditiva, trata
automaticamente tendência e sazonalidades semanal e anual e permite declarar
datas de jogos e feriados como eventos especiais, estimando diretamente seu
efeito médio sobre a série. A comparação é informativa porque os dois modelos
falham por motivos opostos: o primeiro ignora a estrutura temporal se ela não
for codificada em variáveis; o segundo acomoda mal um grande número de
regressores contínuos.

A avaliação segue validação temporal (\emph{out-of-time}): treino em um
intervalo histórico e teste no período subsequente, sem embaralhamento
aleatório, que violaria a dependência temporal da série. As métricas são MAE e
RMSE para a formulação de regressão e AUC e precisão@$k$ para a classificação
de picos, esta última por corresponder ao uso pretendido, em que apenas os dias
de maior risco estimado receberiam alocação adicional de efetivo. Como
referência mínima, os modelos são confrontados com dois \emph{baselines}: a
média histórica por dia da semana e mês, e um modelo sazonal ingênuo que repete
o valor observado sete dias antes. A implementação prevista utiliza Python, com
\emph{scikit-learn}, \emph{xgboost}, \emph{prophet} e PostgreSQL.

\section{Cronograma}

A Tabela~\ref{tab:cronograma} distribui as atividades ao longo das doze semanas
de execução, organizadas em seis quinzenas. A sobreposição entre atividades é
intencional: a obtenção das bases oficiais é o principal gargalo do projeto e
corre em paralelo com a modelagem do banco, e o treinamento começa sobre
subconjuntos já consolidados, sem aguardar a carga completa.

\begin{table}[ht]
\centering
\caption{Cronograma de execução (quinzenas, de 26/08 a 18/11)}
\label{tab:cronograma}
\small
\begin{tabular}{@{}lcccccc@{}}
\toprule
\textbf{Atividade} & \textbf{Q1} & \textbf{Q2} & \textbf{Q3} & \textbf{Q4} &
\textbf{Q5} & \textbf{Q6} \\
\midrule
Revisão bibliográfica e consolidação do escopo & X & & & & & \\
Levantamento e obtenção das bases públicas & X & X & & & & \\
Modelagem do banco relacional e carga dos dados & & X & X & & & \\
Construção da tabela analítica e engenharia de atributos & & & X & X & & \\
Definição do critério de pico e do protocolo de validação & & & X & & & \\
Treinamento e ajuste dos modelos (XGBoost e Prophet) & & & & X & X & \\
Avaliação comparativa e análise de variáveis & & & & & X & X \\
Redação do relatório final e apresentação & & & & & & X \\
\bottomrule
\end{tabular}
\\[4pt]
\footnotesize
Q1: 26/08--08/09 \quad Q2: 09/09--22/09 \quad Q3: 23/09--06/10 \quad
Q4: 07/10--20/10 \quad Q5: 21/10--03/11 \quad Q6: 04/11--18/11
\end{table}

\begin{thebibliography}{9}

\bibitem{cohen1979} Cohen, L.~E. and Felson, M. (1979) ``Social change and
crime rate trends: a routine activity approach''. \emph{American Sociological
Review}, 44(4), pp.~588--608.

\bibitem{kurland2014} Kurland, J., Johnson, S.~D. and Tilley, N. (2014)
``Offenses around stadiums: a natural experiment on crime attraction and
generation''. \emph{Journal of Research in Crime and Delinquency}, 51(1),
pp.~5--28.

\bibitem{ristea2018} Ristea, A., Kurland, J., Resch, B., Leitner, M. and
Langford, C. (2018) ``Estimating the spatial distribution of crime events
around a football stadium from georeferenced tweets''. \emph{ISPRS
International Journal of Geo-Information}, 7(2), 43.

\bibitem{ristea2020} Ristea, A., Al Boni, M., Resch, B., Gerber, M.~S. and
Leitner, M. (2020) ``Spatial crime distribution and prediction for sporting
events using social media''. \emph{International Journal of Geographical
Information Science}, 34(9), pp.~1708--1739.

\end{thebibliography}

\end{document}